%%% SECTION 5: ADAPTED 21 CFR PART 50 FOR PHYSICAL AI HUMAN SUBJECTS PROTECTION %%%

[This section should be approximately 8-10 pages. The major theme is that
Physical AI Adaption: 21 CFR Part 50 \cite{cfr50-adapt} provides comprehensive
human subjects protection specifically designed for Physical AI oncology trials,
building on the existing federal regulation \cite{cfr-part-50}. By adapting
a well-established regulation rather than proposing new legislation, this work
adds immediate credibility and practical applicability to Physical AI
implementations. A critical theme is that control has shifted towards the
patient's side, with more patient rights, including new protections specific
to robotic and AI-assisted medical procedures.]

[PRIMARY SOURCE: national-platform/21cfr50\_adapt/ directory containing three
chunked files: 01\_preamble\_scope\_definitions\_consent.tex (lines 1-438),
02\_irb\_review\_pediatric.tex (lines 439-674), and
03\_additional\_safeguards\_closing.tex (lines 675-747). All three files must
be used together to build this section.]

\subsection{Scope and Physical AI Definitions}
\label{subsec:cfr50-scope}

[Build from national-platform/21cfr50\_adapt/01\_preamble\_scope\_definitions\_consent.tex,
Subpart A (General Provisions). Cover the scope of the adapted regulation and
all Physical AI definitions. The original Part 50 contains 17 Physical AI-specific
definitions that standardize terminology for robotic and AI-assisted clinical
procedures. These definitions are foundational to the entire National Platform.]

[Include key definitions such as Physical AI system, Physical AI adverse event,
Physical AI operator, robotic procedure, autonomous mode, supervised mode,
Physical AI safety matrix, and others. Cross-reference with definitions in the
ICH adaptation (Section~\ref{sec:ich-e6r3}) and the CFR 312 adaptation
(Section~\ref{sec:cfr312}).]

\subsection{Informed Consent for Physical AI Procedures}
\label{subsec:cfr50-consent}

[Build from national-platform/21cfr50\_adapt/01\_preamble\_scope\_definitions\_consent.tex,
Subpart B (Informed Consent of Human Subjects). Cover the adapted general
requirements for informed consent, Physical AI-specific consent elements,
exceptions to informed consent requirements, and documentation requirements.]

[MAJOR THEME: Informed consent for Physical AI procedures is more comprehensive
than traditional consent because patients must understand their interactions
with robotic systems, the AI decision-making processes involved in their care,
the safety monitoring systems in place, and their rights to refuse or pause
robotic procedures at any time. This represents a fundamental shift of control
towards the patient. Reference Patient Instructions: Physical AI Trials
\cite{patient-instructions-paper} as the patient-facing implementation of these
consent requirements.]

[Detail the specific additional consent elements required for Physical AI trials,
such as: description of each robot type and its role, explanation of AI
decision-making in treatment planning, patient's right to request human-only
care at any point, data collection and privacy protections for robot-generated
data, and emergency stop procedures.]

\subsection{Physical AI System Safety Requirements}
\label{subsec:cfr50-safety}

[Build from national-platform/21cfr50\_adapt/02\_irb\_review\_pediatric.tex,
Subpart C (Additional Protections for Subjects in Physical AI Clinical
Investigations). This is a new subpart created specifically for Physical AI
trials. Cover:]

\subsubsection{Pre-Procedure Safety Matrix}
[Detail the safety matrix requirements that must be completed before any
Physical AI procedure begins, including robot system verification, sensor
calibration checks, emergency stop verification, and patient identification
confirmation.]

\subsubsection{Runtime Safety Monitoring}
[Cover continuous monitoring requirements during Physical AI procedures,
including real-time sensor data analysis, deviation detection thresholds,
automatic pause triggers, and human oversight requirements.]

\subsubsection{Post-Procedure Requirements}
[Cover post-procedure documentation, patient follow-up specific to Physical AI
interactions, and adverse event reporting for robot-related incidents.]

\subsubsection{Task-Order Lifecycle}
[Cover the lifecycle management of Physical AI task orders from initiation
through completion, including authorization requirements, execution monitoring,
and completion verification.]

\subsubsection{Forbidden Operations}
[List and explain the operations that Physical AI systems are explicitly
prohibited from performing without human authorization, including certain
surgical decisions, medication administration changes, and emergency
treatment modifications.]

\subsection{IRB Review of Physical AI Investigations}
\label{subsec:cfr50-irb}

[Continue building from national-platform/21cfr50\_adapt/02\_irb\_review\_pediatric.tex.
Cover the specific requirements for IRB review of Physical AI trials, including:
evaluation criteria for Physical AI systems, ongoing monitoring requirements,
protocol modification triggers specific to Physical AI performance, and
criteria for suspending Physical AI operations. Reference the NCI CIRB
\cite{nci-cirb} as a model for centralized Physical AI review.]

\subsection{Ongoing Consent and Subject Notification}
\label{subsec:cfr50-ongoing}

[Build from national-platform/21cfr50\_adapt/02\_irb\_review\_pediatric.tex on
ongoing consent requirements. Cover how participants must be continuously
informed about Physical AI system updates, performance changes, new robot
deployments, and any safety-relevant modifications to the Physical AI systems
they interact with.]

[MAJOR THEME: Unlike traditional trials where informed consent is primarily a
one-time event, Physical AI trials require ongoing consent processes because
the technology may be updated, new robots may be introduced, and AI algorithms
may be refined during the trial. This ongoing consent process gives patients
continuous agency over their participation.]

\subsection{Data Protection and Privacy}
\label{subsec:cfr50-data}

[Build from national-platform/21cfr50\_adapt/02\_irb\_review\_pediatric.tex on
data protection provisions. Cover protections for the additional categories of
data generated by Physical AI systems: robot telemetry, video recordings of
procedures, sensor data from patient-robot interactions, and AI decision logs.
Reference HIPAA \cite{hipaa} and federated learning \cite{fl-paper,
federated-learning-healthcare} as complementary privacy frameworks.]

\subsection{Physical AI System Classification}
\label{subsec:cfr50-classification}

[Build from national-platform/21cfr50\_adapt/02\_irb\_review\_pediatric.tex on
system classification and regulatory pathways. Cover the classification
framework for Physical AI systems by risk level, the corresponding regulatory
pathway for each classification tier, and the criteria for determining which
tier applies. Connect to the USL scoring framework in
Section~\ref{sec:psl-usl}.]

\subsection{Pediatric Protections in Physical AI Trials}
\label{subsec:cfr50-pediatric}

[Build from national-platform/21cfr50\_adapt/02\_irb\_review\_pediatric.tex,
Subpart D (Additional Safeguards for Children). Cover how existing pediatric
protections are adapted for Physical AI contexts, including modified consent
and assent procedures for minors interacting with robots, additional safety
requirements for pediatric Physical AI procedures, and IRB duties specific to
pediatric Physical AI investigations. Note that pediatric oncology is a
particularly important application where Physical AI may reduce treatment
burden on young patients.]

\subsection{Glossary of Physical AI Terms}
\label{subsec:cfr50-glossary}

[Build from national-platform/21cfr50\_adapt/03\_additional\_safeguards\_closing.tex,
Glossary section. Include the supplementary Physical AI terms that complement
the regulatory definitions. Cross-reference with the ICH glossary
(Section~\ref{subsec:ich-glossary}) to ensure consistency across all adapted
standards.]
